% Parte II e III do PDF completo da RSL 2: execução e resultados.
% Gerado a partir dos mesmos arquivos da monografia. Compilar na pasta latex/:
%   pdflatex rsl/rsl2-execucao-resultados && bibtex rsl2-execucao-resultados && pdflatex ... (x2)
\input{fixos/pacotes}
\input{fixos/comandos}
\usepackage{fixos/customizacoes}
\begin{document}
\frenchspacing\openany
\pagestyle{simple}
\begin{center}
{\Large\bfseries RSL 2: Assistentes e agentes de IA no desenvolvimento de software}\\[6pt]
{\large Execução e resultados}\\[12pt]
Matheus Perillo (R1) e orientador (R2) · UnB/FCTE · \today
\end{center}
\vspace{12pt}
\noindent Este documento continua o protocolo da revisão (parte I, nas páginas anteriores). A parte II descreve a execução: a seleção dos estudos, inclusive o \textit{snowballing}, e a avaliação de qualidade. A parte III traz os resultados: o fluxo de seleção, as características dos estudos, a qualidade e a síntese de cada questão de pesquisa. Os textos e as tabelas são os mesmos do capítulo de resultados da monografia, gerados dos mesmos arquivos. A seleção da etapa 2, a avaliação de qualidade, a extração e a síntese foram feitas pela IA como terceira opinião, com regras aceitas por R1. R2 ainda as confirma, e os resultados são provisórios.

\chapter{Execução da revisão}
\section{Seleção dos estudos}
\label{sec:metodologia-selecao}
{\let\subsubsection\subsection \input{editaveis/rsl-execucao-selecao}}
{\let\subsection\section \input{editaveis/rsl-qualidade}}
\section{Síntese}
\label{sec:metodologia-sintese}
\input{editaveis/rsl-sintese-regras}

\chapter{Resultados}
\label{sec:metodologia-rsl}
\phantomsection\label{sec:metodologia-requisitos}
\input{editaveis/gerado/qualidade}
\input{editaveis/resultados-rsl2}

\bibliography{bibliografia}
\appendix
\input{editaveis/gerado/apendice-estudos-rsl2}
\end{document}
